\section{Introduction}
\label{sec:introduction}

\subsection{Background and Motivation}
\label{sec:background_and_motivation}

Driven by electrification and increased hydrogen demand in industry, it is expected that the number of large-scale converters in the system will significantly increase~\cite{elia_2022}. It is well established that large-scale converters cause harmonic pollution, however, the corresponding harmonic frequencies depends on the chosen technology, e.g., a 24-pulse converter under balanced loading conditions mainly causes harmonic distortion at harmonic numbers: 23, 25, 47, 49, etc. The choice of converter should be aligned with the harmonic network impedance to minimize its impact. 

A wide variety of methods exist to measure the harmonic impedance~\cite{robert_1997}. These methods can be categorized into two types: 1) invasive methods, where the harmonic network impedance is extracted from the variations of voltage and current caused by injecting harmonic current; and 2) non-invasive methods, where the harmonic network impedance is extracted from the natural variation of voltage and current. Additionally, these classes can be further divided into a) steady-state methods, and b) transient methods.

The method presented in the paper falls in the transient invasive category, but can be used in both an invasive and non-invasive context. It aims at improving the accuracy of the harmonic network impedance assessment through the use of a novel technique called the Taylor-Fourier transform~\cite{platas_garza_2011}. Given the higher accuracy of the Taylor-Fourier transform in a transient context compared to classic discrete Fourier transform, it becomes possible to accurately assess the harmonic impedance for signals with a lower harmonic content. Typically, for the classic discrete Fourier transform, the relative amplitude of the harmonic current phasor~$\CurrentAmplitude{h}$ should be at least 1\,\% of the fundamental~$\CurrentAmplitude{1}$~\cite{xi_2002}. Through the Taylor-Fourier transform, this factor can be decreased by an order of magnitude, enabling accurate assessment of harmonic impedance at higher harmonic numbers.

\subsection{Literature Review}
\label{sec:literature_review}

Generally, harmonic estimation techniques assume the considered signal to be periodic, implying a constant fundamental frequency and Fourier coefficients, i.e., amplitude and phase, over the entire observation window. The Fourier coefficients are estimated through the discrete Fourier transform, or its `fast' alternative: fast Fourier transform. In practice, signal rarely adhere to the periodic property. For non-periodic signals, the discrete Fourier transform can at best estimate the coefficients by the best average value over the whole observation window, truncating the fluctuation one actually wants to capture.

Two of the main reasons for this are: 1) spectral leakage, i.e., errors produced when the fundamental frequency of the signal fluctuates; and 2) harmonic interference, i.e., errors due to the superposition of wide spectral harmonic components~\cite{liguori_2004}. Another well-known problem of the classic Fourier transform is its lack of time-frequency resolution. Short-time Fourier transform~\cite{gabor_1946,kiymika_2005} tried to solve it but it still presents a constant frequency and time resolution. 

Wavelet theory~\cite{daubechies_1992,vetterli_1995} was proposed to improve resolution in the time-frequency plane, providing good signal compression. However, wavelet transform lacks physical meaning for its coefficients. Furthermore, in general, wavelet transform cannot guarantee a lower reconstruction error than that of the classic discrete Fourier transform. 

Alternatively, harmonics are estimated using an algorithm using recursive approximations~\cite{yang_2005}. Additionally, \cite{tomic_2007} presents an adaptive filtering technique for extracting harmonic components by using a feedback frequency estimation to adapt the filter parameters. However, this method uses non-flat-frequency responses, which are not appropriate for estimating dynamic harmonics; it has delays between the stages; and is only valid for offline applications. 

\subsection{Contributions and Outline}
\label{sec:contributions_and_outline}

The Taylor-Fourier transform relaxes the constant restriction imposed on the Fourier coefficients of the signal model. Each harmonic is modulated by a time function, which is given by a Taylor polynomial of a given order greater than zero. In this way, the coefficients of the Taylor-Fourier transform include, in addition to the instantaneous value of each harmonic fluctuation, its first-order derivatives as in a Taylor approximation. Concretely, the contributions of the paper include:
\begin{enumerate}
    \item a method based on the first-order derivative of the transient voltage and current signals to estimate the harmonic impedance using the Taylor-Fourier transform, and
    \item a case study on transient signals measured during the inrush of a power transformer at an industrial power system, proving the effectiveness of the method.
\end{enumerate}
Furthermore, the Taylor-Fourier transform approach is implemented in Julia, and made available as an open-source package: \textsc{TaylorFourierTransform.jl}\footnote{Available at: \href{https://github.com/timmyfaraday/TaylorFourierTransform.jl}{https://github.com/timmyfaraday/TaylorFourierTransform.jl}}.

The remainder of the paper is organized as follows: Section~\ref{sec:mathematical_framework} presents the mathematical framework, including theoretical background on the Taylor-Fourier Transform. Section~\ref{sec:case_study} shows numerical results for the case study on the transient signals measured during the inrush of a power transformer. Section~\ref{sec:conclusions} concludes the paper.